\documentclass{article}
\usepackage{geometry}
\geometry{margin=1.5cm, vmargin={0pt,1cm}}
\setlength{\topmargin}{-1cm}
\setlength{\headheight}{12.5pt}
\setlength{\paperheight}{29.7cm}
\setlength{\textheight}{25.3cm}

% useful packages.
\usepackage[UTF8]{ctex}
\usepackage{fancyhdr}
\usepackage{layout}
\usepackage{luacode}

% import common macros
% % % % % % % % % % % % % % % % % % % % % % % % % % % % % % % %
% Common Macros for LaTeX
% % % % % % % % % % % % % % % % % % % % % % % % % % % % % % % %

% Useful packages
\usepackage{amsfonts}
\usepackage{amsmath}
\usepackage{amssymb}
\usepackage{amsthm}
\usepackage{enumerate}
\usepackage{graphicx}
\usepackage{multicol}
\usepackage[ruled,linesnumbered]{algorithm2e}

% Math operators and common symbols
\newcommand{\dif}{\mathrm{d}}
\newcommand{\innerProd}[1]{\left\langle #1 \right\rangle}
\newcommand{\difFrac}[2]{\frac{\dif #1}{\dif #2}}
\newcommand{\pdfFrac}[2]{\frac{\partial #1}{\partial #2}}
\newcommand{\T}{\mathrm{T}}
\newcommand{\norm}[1]{\left\| #1 \right\|}
\newcommand{\abs}[1]{\left| #1 \right|}

% Vector and Matrix shortcuts
\newcommand{\vecTwo}[2]{
  \begin{bmatrix}
    #1 \\
    #2
\end{bmatrix}}
\newcommand{\vecThree}[3]{
  \begin{bmatrix}
    #1 \\
    #2 \\
    #3
\end{bmatrix}}
\newcommand{\vecTwoH}[2]{
  \begin{bmatrix}
    #1 & #2
  \end{bmatrix}
}
\newcommand{\vecThreeH}[3]{
  \begin{bmatrix}
    #1 & #2 & #3
  \end{bmatrix}
}

% Common fields
\newcommand{\R}{\mathbb{R}}
\newcommand{\C}{\mathbb{C}}
\newcommand{\Z}{\mathbb{Z}}
\newcommand{\N}{\mathbb{N}}

% Solutions and proofs
\newenvironment{solution}{
\begin{proof}[解]}{
\end{proof}}

% Utility to get environment variables (requires lualatex)
\newcommand{\getEnv}[2]{\directlua{tex.print(os.getenv("#1") or "#2")}}


\usepackage{listings}
\usepackage{xcolor}

% Configure listings for Python code highlighting
\lstset{
  language=Python,
  basicstyle=\ttfamily\small,
  keywordstyle=\bfseries\color{blue},
  commentstyle=\color{gray},
  stringstyle=\color{red},
  breaklines=true,
  showstringspaces=false,
  tabsize=4,
  frame=single,
  framesep=5pt,
  rulecolor=\color{gray},
  backgroundcolor=\color{white},
  captionpos=b,
  numberstyle=\tiny\color{gray},
  numbers=left,
  stepnumber=5,
}

\lstnewenvironment{plaintext}[1][]{
  \lstset{
    language=none,
    basicstyle=\ttfamily,
    keywordstyle=,
    commentstyle=,
    stringstyle=,
    numbers=none,
    frame=none,
    backgroundcolor=\color{white},
    #1                % 允许用户自行覆盖默认设置
  }
}{}

% Name and uID
\newcommand{\name}{\getEnv{NAME}{NAME}}
\newcommand{\uid}{\getEnv{UID}{UID}}
\newcommand{\course}{数值代数}
\newcommand{\hwNum}{12}

\begin{document}

\pagestyle{fancy}
\fancyhead{}
\lhead{\name (\uid)}
\chead{\course \#\hwNum}
\rhead{\today}

\section{题目 1}

分别应用幂法于矩阵
\[
  A =
  \begin{pmatrix} \lambda & 1 \\ 0 & \lambda
  \end{pmatrix} \quad \text{和} \quad B =
  \begin{pmatrix} \lambda & 1 \\ 0 & -\lambda
  \end{pmatrix} \quad (\lambda > 0)
\]
并考虑所得序列的特性, 求出特征值和特征向量.

\begin{solution}

  \textbf{矩阵 $A$ 的幂法分析.}

  取初始向量 $\boldsymbol{v}^{(0)} = (1, \delta)^\T$ ($\delta \neq 0$).
  注意到 $A = \lambda I + N$，其中
  \[
    N =
    \begin{pmatrix} 0 & 1 \\ 0 & 0
    \end{pmatrix}, \quad N^2 = 0.
  \]
  由二项式展开（因 $N^2 = 0$，高次项消失）：
  \[
    A^k = (\lambda I + N)^k = \lambda^k I + k\lambda^{k-1}N
    = \lambda^k
    \begin{pmatrix} 1 & k/\lambda \\ 0 & 1
    \end{pmatrix}.
  \]
  对初始向量 $\boldsymbol{v}^{(0)} = (1, \delta)^\T$：
  \[
    A^k \boldsymbol{v}^{(0)} = \lambda^k
    \begin{pmatrix} 1 + k\delta/\lambda \\ \delta
    \end{pmatrix}.
  \]

  在幂法中，按首个分量归一化。
  设归一化后的向量 $\boldsymbol{v}^{(k)} = (1, c_k)^\T$，则迭代关系为
  \[
    \boldsymbol{y}^{(k+1)} = A\boldsymbol{v}^{(k)}
    =
    \begin{pmatrix} \lambda + c_k \\ \lambda c_k
    \end{pmatrix},
    \quad \mu_{k+1} = \lambda + c_k,
    \quad c_{k+1} = \frac{\lambda c_k}{\lambda + c_k}.
  \]
  由数学归纳法可得
  \[
    c_k = \frac{\lambda\delta}{\lambda + k\delta} \xrightarrow{k\to\infty} 0.
  \]

  因此：
  \begin{itemize}
    \item 特征值估计：$\mu_k = \lambda + c_{k-1} = \lambda +
      \dfrac{\lambda\delta}{\lambda + (k-1)\delta} \to \lambda$，
      误差为 $O(1/k)$；
    \item 特征向量：$\boldsymbol{v}^{(k)} = (1, c_k)^\T \to (1, 0)^\T$。
  \end{itemize}

  \textbf{序列特性：}
  $A$ 的特征值 $\lambda$ 的代数重数为 2，但几何重数为 1（亏损特征值），
  导致幂法收敛速度为 $O(1/k)$，慢于正常的几何收敛速度。

  \medskip

  \textbf{矩阵 $B$ 的幂法分析.}

  注意到
  \[
    B^2 =
    \begin{pmatrix} \lambda & 1 \\ 0 & -\lambda
    \end{pmatrix}^2
    =
    \begin{pmatrix} \lambda^2 & 0 \\ 0 & \lambda^2
    \end{pmatrix}
    = \lambda^2 I.
  \]
  因此 $B^{2m} = \lambda^{2m}I$，$B^{2m+1} = \lambda^{2m}B$。

  取初始向量 $\boldsymbol{v}^{(0)} = (1, \delta)^\T$ ($\delta \neq 0$)。
  按首个分量归一化：
  \begin{align*}
    \boldsymbol{y}^{(1)} &= B\boldsymbol{v}^{(0)}
    =
    \begin{pmatrix} \lambda + \delta \\ -\lambda\delta
    \end{pmatrix},
    &\mu_1 &= \lambda + \delta,
    &\boldsymbol{v}^{(1)} &= \left(1,\;
    \frac{-\lambda\delta}{\lambda+\delta}\right)^\T. \\[6pt]
    \boldsymbol{y}^{(2)} &= B\boldsymbol{v}^{(1)}
    =
    \begin{pmatrix} \lambda^2/(\lambda+\delta)
      \\ \lambda^2\delta/(\lambda+\delta)
    \end{pmatrix},
    &\mu_2 &= \frac{\lambda^2}{\lambda+\delta},
    &\boldsymbol{v}^{(2)} &= (1,\; \delta)^\T = \boldsymbol{v}^{(0)}.
  \end{align*}

  因此序列以周期 2 振荡：
  $\boldsymbol{v}^{(0)}, \boldsymbol{v}^{(1)}, \boldsymbol{v}^{(0)},
  \boldsymbol{v}^{(1)}, \ldots$，不收敛。

  \textbf{特征值：}两个振荡的 $\mu$ 值之积为
  \[
    \mu_1 \cdot \mu_2 = (\lambda+\delta) \cdot
    \frac{\lambda^2}{\lambda+\delta} = \lambda^2,
  \]
  故 $|\lambda| = \sqrt{\mu_1\mu_2}$。
  由 $\lambda > 0$ 知 $\lambda = \sqrt{\mu_1\mu_2}$。
  又由 $\mathrm{tr}\, B = 0$ 知两特征值之和为 0，故另一特征值为 $-\lambda$。

  \textbf{特征向量：}
  \begin{itemize}
    \item 对 $\lambda_1 = \lambda$：$(B - \lambda I)\boldsymbol{v} =
      \boldsymbol{0}$，
      即 $
      \begin{pmatrix} 0 & 1 \\ 0 & -2\lambda
      \end{pmatrix}\boldsymbol{v} = \boldsymbol{0}$，
      解得 $\boldsymbol{v}_1 = (1,\; 0)^\T$；
    \item 对 $\lambda_2 = -\lambda$：$(B + \lambda I)\boldsymbol{v} =
      \boldsymbol{0}$，
      即 $
      \begin{pmatrix} 2\lambda & 1 \\ 0 & 0
      \end{pmatrix}\boldsymbol{v} = \boldsymbol{0}$，
      解得 $\boldsymbol{v}_2 = (1,\; -2\lambda)^\T$。
  \end{itemize}

  \textbf{序列特性：}
  $B$ 的两个特征值 $\lambda$ 和 $-\lambda$ 具有相同的模 $|\lambda|$，
  幂法无法收敛至单一特征值，序列以周期 2 振荡。
  但通过 $\sqrt{\mu_1\mu_2}$ 仍可恢复 $|\lambda|$。

\end{solution}

\section{题目 2}

求多项式方程的模最大根.

(1) 用你所熟悉的计算机语言编制利用幂法求多项式方程
\[
  f(x) = x^n + \alpha_{n-1}x^{n-1} + \dots + \alpha_1x + \alpha_0 = 0
\]
的模最大根的通用子程序.

(2) 利用你所编制的子程序求下列各高次方程的模最大根:

(i) $x^3 + x^2 - 5x + 3 = 0$;

(ii) $x^3 - 3x - 1 = 0$;

(iii) $x^8 + 101x^7 + 208.01x^6 + 10891.01x^5 + 9802.08x^4 +
79108.9x^3 - 99902x^2 + 790x - 1000 = 0$.

\begin{solution}

  \textbf{方法.}
  设多项式方程
  \[
    f(x) = x^n + a_{n-1}x^{n-1} + \dots + a_1 x + a_0 = 0,
  \]
  则其根恰为友矩阵 (companion matrix)
  \[
    C =
    \begin{pmatrix}
      0 & 0 & \cdots & 0 & -a_0 \\
      1 & 0 & \cdots & 0 & -a_1 \\
      0 & 1 & \cdots & 0 & -a_2 \\
      \vdots & \vdots & \ddots & \vdots & \vdots \\
      0 & 0 & \cdots & 1 & -a_{n-1}
    \end{pmatrix}
  \]
  的特征值。因此求模最大根等价于求 $C$ 的模最大特征值，可用幂法求解。

  幂法迭代格式：取初始向量 $\boldsymbol{v}^{(0)}$，对 $k = 0, 1, 2, \dots$ 计算
  \[
    \boldsymbol{y}^{(k+1)} = C\boldsymbol{v}^{(k)}, \quad
    \mu_{k+1} = y_j^{(k+1)} \quad (j = \arg\max_i |y_i^{(k+1)}|), \quad
    \boldsymbol{v}^{(k+1)} = \boldsymbol{y}^{(k+1)} / \mu_{k+1}.
  \]
  当 $|\mu_k|$ 收敛后，通过 Rayleigh 商
  $\lambda = \boldsymbol{v}^\T C \boldsymbol{v} / (\boldsymbol{v}^\T
  \boldsymbol{v})$
  确定特征值的符号。

  代码与运行结果如下：

  \lstinputlisting[language=Python, caption={幂法求多项式模最大根}]{exercise.py}
  \lstinputlisting[language={}, caption={运行结果}]{output.txt}

  \textbf{结果分析.}
  \begin{enumerate}[(i)]
    \item $x^3 + x^2 - 5x + 3 = (x-1)^2(x+3) = 0$，
      模最大根为 $x = -3$，幂法经 28 次迭代收敛至 $-3.0000000000$。
      由于 $\lambda_1/\lambda_2 = -3/1 = -3$，收敛比为 $1/3$。
    \item $x^3 - 3x - 1 = 0$，
      模最大根为 $x \approx 1.8793852416$，
      经 120 次迭代收敛。
      收敛比为 $|\lambda_2/\lambda_1| = 1.5321/1.8794 \approx 0.8154$，收敛较慢。
    \item 八次多项式的模最大根为 $x = -100$，
      经 15 次迭代即收敛。
      该多项式的 8 个根为 $-100,\; \pm 10i,\; -1\pm 3i,\; 1,\; \pm 0.1i$，
      收敛比为 $10/100 = 0.1$，收敛很快。
  \end{enumerate}

\end{solution}

\end{document}
